%% ma: L12-B05-0005
%% lop: 12
%% chuong: 1
%% bai: 5
%% dang: 12.5.1
%% loai: TN4
%% muc: VD
%% dap_an: "A"
%% nguon: "(NBV)-ĐỀ SỐ 4-MỖI NGÀY 1 ĐỀ THI 2025.pdf (D04, MỖI NGÀY 1 ĐỀ - Đề số 4), Phần I Câu 4"
%% kien_thuc: "Bài 5 (ứng dụng đạo hàm vào tối ưu thực tiễn; biến nguyên)"
%% trang_thai: chinh-thuc
%% ghi_chu: "Cùng dạng 12.5.1 với L12-B05-0001 (tối ưu lợi nhuận, biến nguyên) nhưng hàm lợi nhuận cho sẵn và nghiệm nguyên. Đáp án gốc A đúng (sympy: x=24)."
\begin{ex}
Giả sử lợi nhuận của một phân xưởng khi sản xuất và bán ra $x$ chiếc xe đạp mỗi tuần được tính theo công thức $L(x)=-x^3+3x^2+1584x-200$ (nghìn đồng) ($0<x\le40$). Mỗi tuần phân xưởng đó sản xuất và bán ra bao nhiêu chiếc xe đạp thì có lợi nhuận cao nhất?
\choice{\True $24$}{$13$}{$10$}{$22$}
\loigiai{$L'(x)=-3x^2+6x+1584=0\Leftrightarrow x=24$ (nhận) hoặc $x=-22$ (loại). $L'(x)>0$ trên $(0;24)$ và $L'(x)<0$ trên $(24;40]$ nên $L$ đạt giá trị lớn nhất tại $x=24$.}
\end{ex}
